\documentclass[../main.tex]{subfiles}

\begin{document}

\chapter{ Week 1 }

代靖涵 25120222201319

\begin{exercise}{}{9-2-1}
     设 \(  \left\{ f_{k} \right\}  \)为实值函数列且 \(  g\left( x \right)= \sup _{k\in \mathbb{N} }f_{k}\left( x \right)    \), \(  h\left( x \right)= \inf _{k\in \mathbb{N} }f_{k}\left( x \right)    \). 证明: \(  \forall a\in \mathbb{R}   \)
    \begin{enumerate}
        \item \[
        \left\{ x:g\left( x \right)> a  \right\}= \bigcup _{k= 1}^{\infty}\left\{ x:f_{k}\left( x \right)> a  \right\},
        \]
        \item \[
        \left\{ x: g\left( x \right)\le a  \right\}= \bigcap_{k= 1}^{\infty}\left\{ x:f_{k}\left( x \right)\le a  \right\} ,
        \]
        \item \[
        \left\{ x:h\left( x \right)< a  \right\}= \bigcup _{k= 1}^{\infty}\left\{ x:f_{k}\left( x \right)< a  \right\},
        \]
        \item \[
        \left\{ x:h\left( x \right)\ge a  \right\}= \bigcap_{k= 1}^{\infty}\left\{ x:f_{k}\left( x \right)\ge a  \right\} ,
        \]
    \end{enumerate}
        
\end{exercise}

\begin{proof}
    \begin{enumerate}
        \item 给定 \(  x  \), \(  g\left( x \right)> a   \)当且仅当 \(  a  \)不是 \(  \left\{ f_{k}\left( x \right)  \right\}  \)的上界,这当且仅当存在 \(  k\in \mathbb{N}   \),使得 \(  f_{k}\left( x \right)> a   \).即 存在 \(  k \in \mathbb{N}   \), 使得 \(  x \in \left\{ x:f_{k}\left( x \right)> a  \right\}  \),也就是说 \[
        x \in \bigcup _{k= 1}^{\infty}\left\{ x:f_{k}\left( x \right)> a  \right\}
        \]因此 \[
        \left\{ x:g\left( x \right)> a  \right\}= \bigcup _{k= 1}^{\infty}\left\{ x:f_{k}\left( x \right)> a  \right\}
        \]  
        \item  \(  g\left( x \right)\le a   \), 当且仅当 \(  a  \)是 \(  \left\{ f_{k}\left( x \right)  \right\}  \)的上界; 当且仅当对于任意的   \(  k\in \mathbb{N}   \), \(  f_{k}\left( x \right)\le a   \); 当且仅当 \[
        x \in \bigcap_{k= 1}^{\infty}\left\{ x:f_{k}\left( x \right)\le a  \right\} 
        \]
        \item 注意到 \[
        h\left( x \right)=  -\sup _{k\in \mathbb{N} }(-f_{k}\left( x \right) ) 
        \] \(  h\left( x \right)< a   \)当且仅当 \(  \sup _{k\in \mathbb{N} }(-f_{k}\left( x \right) )> -a  \) . 将1.的结果应用于实值函数列 \(  \left\{ -f_{k}\left( x \right)  \right\}  \)以及值 \( - a  \), 得到 \[
        \begin{aligned}
        \left\{ x:h\left( x \right)< a  \right\}= \left\{ x: \sup _{k\in \mathbb{N} }\left( -f_{k}\left( x \right) \right) > -a  \right\}&= \bigcup _{k= 1}^{\infty} \left\{ x: -f_{k}\left( x \right) > -a \right\}\\ 
         &= \bigcup _{k= 1}^{\infty}\left\{ x:  f_{k}\left( x \right)< a \right\}
        \end{aligned}
        \]   
        \item 类似地, \(  h\left( x \right)\ge a   \)当且仅当 \(  \sup _{k\in \mathbb{N} }(-f_{k}\left( x \right) )\le -a  \). 将2.的结果应用与实值函数列 \(  \left\{ -f_{k}\left( x \right)  \right\}  \)以及值 \(  -a  \), 得到   \[
        \begin{aligned}
        \left\{ x:h\left( x \right)\ge a  \right\}= \left\{ x:\sup _{k\in \mathbb{N} }(-f_{k}\left( x \right) )\le -a \right\}&= \bigcap _{k= 1}^{\infty}\left\{ x:-f_{k}\left( x \right)\le -a  \right\}\\ 
         &= \bigcap _{k= 1}^{\infty}\left\{ x:f_{k}\left( x \right)\ge a  \right\} 
        \end{aligned}
        \]  
    \end{enumerate}
    
\end{proof}

\begin{exercise}{}{9-2-2}
    设 $\{f_m\}$ 为实值函数列. 下列等式是否成立? 若成立请证明, 若不成立请举出反例.
\begin{enumerate}
\item[(i)] $\bigcup_{k=1}^{\infty} \bigcap_{m=1}^{\infty} \{x : f_m(x) \ge \frac{1}{k}\} = \bigcup_{k=1}^{\infty} \bigcap_{m=1}^{\infty} \{x : f_m(x) > \frac{1}{k}\};$
\item[(ii)] $\bigcap_{k=1}^{\infty} \bigcup_{m=1}^{\infty} \{x : f_m(x) \le \frac{1}{k}\} = \bigcap_{k=1}^{\infty} \bigcup_{m=1}^{\infty} \{x : f_m(x) < \frac{1}{k}\}.$
\end{enumerate}
\end{exercise}
\begin{proof}
  \begin{enumerate}
    \item  由于  \[
 \left\{ x:f_{m}\left( x \right)\ge \frac{1 }{k }   \right\}\supseteq  \left\{ x:f_{m}\left( x \right)> \frac{1 }{k }   \right\} 
    \]故 \[
    LHS\supseteq RHS
    \]另一方面, 注意到 \[
    \left\{ x:f_{m}\left( x \right)\ge \frac{1 }{k }   \right\}\subseteq \left\{ x:f_{m}\left( x \right)> \frac{1 }{k+ 1 }   \right\}
    \]我们有 \[
    \bigcup _{k= 1}^{\infty}\bigcap_{m= 1}^{\infty}\left\{ x:f_{m}\left( x \right)\ge \frac{1 }{k }   \right\} \subseteq \bigcup _{k= 2}^{\infty}\bigcap_{m= 1}^{\infty}\left\{ x:f_{m}\left( x \right)> \frac{1 }{k }   \right\}\subseteq \bigcup _{k= 1}^{\infty}\bigcap_{m= 1}^{\infty}\left\{ x:f_{m}\left( x \right)> \frac{1 }{k }   \right\}  
    \]因此两集合相等. 
    \item 由于 \[
    \left\{ x:f_{m}\left( x \right)\le \frac{1 }{k }   \right\}\supseteq \left\{ x:f_{m}\left( x \right)< \frac{1 }{k }   \right\}
    \]故 \[
    LHS\supseteq RHS
    \]另一方面,  \[
    \left\{ x:f_{m}\left( x \right)\le \frac{1 }{k }   \right\}\subseteq \left\{ x:f_{m}\left( x \right)< \frac{1 }{k-1 }   \right\},\quad \forall k\ge 2
    \] \[
   \begin{aligned}
    \bigcap_{k= 1}^{\infty}\bigcup _{m= 1}^{\infty}\left\{ x:f_{m}\left( x \right)\le \frac{1 }{k }   \right\}&\subseteq \bigcap_{k= 2}^{\infty}\bigcup _{m= 1}^{\infty}\left\{ x:f_{m}\left( x \right)\le \frac{1 }{k }   \right\}\\ 
     &\subseteq \bigcap_{k= 2}^{\infty}\bigcup _{m= 1}^{\infty}\left\{ x:f_{m}\left( x \right)< \frac{1 }{k-1 }   \right\} \\ 
      &=  \bigcap_{k= 1}^{\infty}\bigcup _{m= 1}^{\infty}\left\{ x:f_{m}\left( x \right)< \frac{1 }{k }   \right\}    
   \end{aligned}
    \]因此两集合相等.
  \end{enumerate}
  
\end{proof}

\begin{exercise}{}{9-2-3}
    设 $A_1 \subset A_2 \subset \cdots \subset A_n \subset \cdots$, $B_1 \subset B_2 \subset \cdots \subset B_n \subset \cdots$, 则
\[ \left( \bigcup_{n=1}^{\infty} A_n \right) \cap \left( \bigcup_{n=1}^{\infty} B_n \right) = \bigcup_{n=1}^{\infty} (A_n \cap B_n). \]
\end{exercise}
\begin{proof}
    \[
    \begin{aligned}
    x\in LHS&\iff \exists n_1,n_2\in \mathbb{N} , s.t. x\in A_{n_1}\cap B_{n_2} 
    \end{aligned}
    \] 不妨设 \(  n_2\ge n_1  \), 则\[
    x \in A_{n_1}\cap B_{n_2}\implies  x \in A_{n_2}\cap B_{n_2}\subseteq \bigcup _{n = 1}^{\infty}\left( A_{n}\cap B_{n} \right) 
    \] 因此 \[
    LHS \subseteq RHS
    \] 另一方面, 若 \(  x \in \bigcup _{n = 1}^{\infty}\left( A_{n}\cap B_{n} \right)   \), 则存在 \(  n_0\in \mathbb{Z} _{> 0}  \), 使得 \[
    x\in A_{n_0}\cap B_{n_0}
    \]从而 \[
    x \in A_{n_0}\subseteq \bigcup _{n = 1}^{\infty}A_{n},\quad x\in B_{n_0}\subseteq \bigcup _{n = 1}^{\infty}B_{n}
    \]进而 \[
    x \in \left( \bigcup _{n = 1}^{\infty}A_{n} \right)\cup \left( \bigcup  _{n =  1}^{\infty}B_{n}  \right) 
    \]  这表明 \[
    RHS\supseteq LHS
    \]综上, 命题中两集合相等.
\end{proof}

\begin{exercise}{}{}
    设有集合 $A,B,E$ 和 $F$.
\begin{enumerate}
\item[(i)] 若 $A \cup B = E \cup F, A \cap F = \emptyset$ 以及 $B \cap E = \emptyset$, 则
$A = E \text{且} B = F;$
\item[(ii)] 若 $A \cup B = E \cup F$, 令 $A_1 = A \cap E, A_2 = A \cap F$, 则
$A_1 \cup A_2 = A.$
\end{enumerate}
\end{exercise}

 \begin{proof}
 \begin{enumerate}
    \item[(i)]  \[
  E\subseteq  E\cup F= A\cup B
  \] 由于 \(  B\cap E= \varnothing  \), 因此 \[
  E\subseteq  A
  \] 类似地, \[
  F\subseteq E\cup F= A\cup B
  \], 由于 \(  A\cap F= \varnothing  \), 我们有 \[
  F\subseteq B
  \] 由对称性, 通过调换 \(  E  \)和 \(  A  \), \(  B  \)和 \(  F  \)的位置, 可得 \[
  A\subseteq E,\quad B\subseteq F
  \]   因此 \[
  A= E,\quad B= F
  \]
  \item[(ii)]  \[
  A_1\cup A_2= \left( A\cap E \right)\cup \left( A\cap F \right)= A\cap \left( E\cup F \right)= A\cap \left( A\cup B \right)= A    
  \]
 \end{enumerate}
 
 \end{proof}

 \begin{exercise}{}{}
设$\{f_j(x)\}$是定义在$\mathbf{R}^n$上的函数列, 试用点集 $\{x: f_j(x) \ge 1/k\} \quad (j,k = 1,2,\cdots)$ 表示点集$\{x: \varlimsup_{j\to\infty} f_j(x)>0\}$.
\end{exercise}
\begin{proof}
     \[
     \overline{\lim}_{j\to \infty}f_{j}\left( x \right)= \inf _{j\in \mathbb{Z} _{> 0}} \sup _{m\ge j}f_{m}\left( x \right)  
     \] 令\(  g_{j}\left( x \right): = \sup _{m\ge j}f_{m}\left( x \right)    \). 则\[
     \left\{ x: \overline{\lim}_{j\to \infty}f_{j}\left( x \right)> 0  \right\}= \left\{ x: \inf _{j\in \mathbb{Z} _{> 0}}g_{j}\left( x \right)> 0  \right\}
     \] 其中 \[
     \left\{ x:\inf _{j\in \mathbb{Z} _{> 0}}g_{j}\left( x \right)> 0  \right\}= \bigcup _{k= 1}^{\infty}\left\{ x: \inf _{j\in \mathbb{Z} _{> 0}}g_{j}\left( x \right)\ge \frac{1 }{k }   \right\}
     \]由Exercise \ref{ex:9-2-1}可知 \[
     \left\{ x:\inf _{j\in \mathbb{Z} _{> 0}}g_{j}\left( x \right)\ge \frac{1 }{k }   \right\}= \bigcap_{j= 1}^{\infty}\left\{ x:g_{j}\left( x \right)\ge \frac{1 }{k }   \right\} 
     \]因此 \[
     \left\{ x:\overline{\lim}_{j\to \infty}f_{j}\left( x \right)> 0  \right\}= \bigcup _{k= 1}^{\infty}\bigcap_{j= 1}^{\infty}\left\{ x: \sup _{m\ge j}f_{m}\left( x \right)\ge \frac{1 }{k }   \right\} 
     \] 对函数列 \(  \left\{ \sup _{m\ge j}f_{m}\left( x \right)  \right\}_{j}  \)应用 Exercise  \ref{ex:9-2-2} (i), 得到 \[
     \bigcup _{k= 1}^{\infty}\bigcap_{j= 1}^{\infty}\left\{ x: \sup _{m\ge j}f_{m}\left( x \right)\ge \frac{1 }{k }   \right\}= \bigcup _{k= 1}^{\infty}\bigcap_{j= 1}^{\infty}\left\{ x: \sup _{m\ge j}f_{m}\left( x \right)> \frac{1 }{k }   \right\}  
     \]再由 Exercise \ref{ex:9-2-1} 1.的结果,得到 \[
     \left\{ x:\sup _{m\ge j}f_{m}\left( x \right)> \frac{1 }{k }   \right\}= \bigcup _{m= j}^{\infty}\left\{ x:f_{m}\left( x \right)> \frac{1 }{k }   \right\}
     \]最终我们有 \[
     \left\{ x: \overline{\lim}_{j\to \infty}f_{j}\left( x \right)> 0  \right\}= \bigcup _{k= 1}^{\infty}\bigcap_{j= 1}^{\infty}\bigcup _{m= j}^{\infty}\left\{ x:f_{m}\left( x \right)> \frac{1 }{k }   \right\} 
     \]断言 \[
     \bigcup _{k= 1}^{\infty}\bigcap_{j= 1}^{\infty}\bigcup _{m= j}^{\infty}\left\{ x:f_{m}\left( x \right)> \frac{1 }{k }   \right\}=     \bigcup _{k= 1}^{\infty}\bigcap_{j= 1}^{\infty}\bigcup _{m= j}^{\infty}\left\{ x:f_{m}\left( x \right)\ge  \frac{1 }{k }   \right\}
     \]其中 \( LHS  \subseteq   RHS\)由 \(  \left\{ x:f_{m}\left( x \right)> \frac{1 }{k }   \right\}\subseteq \left\{ x:f_{m}\left( x \right)\ge \frac{1 }{k }   \right\}  \)得到.  另一边注意到 \[
    \left\{ x:f_{m}\left( x \right)> \frac{1 }{k+ 1 }   \right\}\supseteq \left\{ x:f_{m}\left( x \right)\ge \frac{1 }{k}   \right\}
     \], 从而 \[
    \begin{aligned}
     \bigcup _{k= 1}^{\infty}\bigcap_{j= 1}^{\infty}\bigcup _{m= j}^{\infty}\left\{ x:f_{m}\left( x \right)\ge  \frac{1 }{k }   \right\}& \subseteq \bigcup _{k= 2}^{\infty}\bigcap_{j= 1}^{\infty}\bigcup _{m= j}^{\infty}\left\{ x:f_{m}\left( x \right)>  \frac{1 }{k }   \right\}  \\ 
      & \subseteq \bigcup _{k= 1}^{\infty}\bigcap_{j= 1}^{\infty}\bigcup _{m= j}^{\infty}\left\{ x:f_{m}\left( x \right)> \frac{1 }{k }   \right\} 
    \end{aligned}
     \]因此\[
     \left\{ x: \overline{\lim}_{j\to \infty}f_{j}\left( x \right)> 0  \right\}= \bigcup _{k= 1}^{\infty}\bigcap_{j= 1}^{\infty}\bigcup _{m= j}^{\infty}\left\{ x:f_{m}\left( x \right)\ge  \frac{1 }{k }   \right\} 
     \]
\end{proof}
\begin{exercise}{}{}
设$\{f_n(x)\}$是定义在$[a,b]$上的函数列, $E\subset [a,b]$且有
\[\lim_{n\to\infty}f_n(x) = \chi_{[a,b]\setminus E}(x), \quad x\in [a,b].\]
若令 $E_n = \{x\in [a,b]: f_n(x)\ge \frac{1}{2}\}$, 试求集合 $\lim_{n\to\infty} E_n$.
\end{exercise}
\begin{proof}
   若 \(  x \in [a,b]\setminus E  \), 则 \[
   \lim_{n\to \infty}f_{n}\left( x \right)= 1
   \]从而 存在 \(  k \in \mathbb{Z} _{> 0}  \), 使得对于任意的 \(  m\ge k  \), 都有 \(  f_{m}\left( x \right)\ge  \frac{1 }{2 }    \), 即 \[
   x \in \bigcup _{k = 1}^{\infty}\bigcap_{m = k}^{\infty}E_{m}=  \liminf_{n\to \infty}E_{n} 
   \]   这表明 \[
   [a,b]\setminus E\subseteq  \liminf_{n\to \infty}E_{n}
   \] 若 \(  x \in E  \), 则 \(  \lim_{n\to \infty}f_{n}\left( x \right)= 0   \), 存在 \(  k \in \mathbb{Z} _{> 0}  \), 使得对于任意的 \(  m\ge k  \), \(  f_{m}\left( x \right)< \frac{1 }{2 }    \)     , 即 \[
   \begin{aligned}
    x\in \bigcup _{k= 1}^{\infty}\bigcap_{m= k}^{\infty}E_{m}^{c}=  \bigcup _{k = 1}^{\infty}\left( \bigcup _{m = k}^{\infty}E_{m} \right)^{c}&= \left( \bigcap_{k= 1}^{\infty}\bigcup _{m= k}^{\infty}E_{m}  \right)^{c}    \\ 
     &=\left(  \limsup_{n\to \infty}E_{n} \right)^{c} 
   \end{aligned}
   \]这表明 \[
   E\subseteq \left( \limsup_{n\to \infty}E_{n} \right)^{c} 
   \]即 \[
   [a,b]\setminus E\supseteq \limsup_{n\to \infty}E_{n}
   \]综上, \[
   \left[ a,b \right]\setminus E \subseteq \liminf_{n\to \infty}E_{n}\subseteq \limsup_{n\to \infty}E_{n}\subseteq \left[ a,b \right]\setminus E 
   \]这表明 \(  \lim_{n\to \infty}E_{n}  \)存在, 且 \[
   \lim_{n\to \infty}E_{n}= \left[ a,b \right]\setminus E 
   \] 
\end{proof}
\begin{exercise}{}{}
设有集合列$\{A_n\}, \{B_n\}$, 试证明
\begin{enumerate}
    \item[(i)] $\varlimsup_{n\to\infty} (A_n \cup B_n) = \left(\varlimsup_{n\to\infty} A_n\right) \cup \left(\varlimsup_{n\to\infty} B_n\right);$
    \item[(ii)] $\varliminf_{n\to\infty} (A_n \cap B_n) = \left(\varliminf_{n\to\infty} A_n\right) \cap \left(\varliminf_{n\to\infty} B_n\right).$
\end{enumerate}
\end{exercise}
\begin{proof}
    \begin{enumerate}
        \item 令 \(  C_{k}= \bigcup _{n = k}^{\infty}A_{n}   \), \(  D_{k}= \bigcup _{n = k}^{\infty}B_{n}  \). 则 \[
    \left( \overline{\lim}_{n\to \infty}A_{n} \right)\cup \left( \overline{\lim}_{n\to \infty}B_{n} \right)=\left(  \bigcap_{k= 1}^{\infty}C_{k} \right) \cup    \left( \bigcap_{k= 1}^{\infty}D_{k}  \right) 
    \]   \[
    \begin{aligned}
    \overline{\lim_{n\to \infty}}\left( A_{n}\cup B_{n} \right)=  \bigcap_{k= 1}^{\infty}\bigcup _{n = k}^{\infty}\left( A_{n}\cup B_{n} \right)&= \bigcap_{k= 1}^{\infty}\left(   \left( \bigcup _{n = k}^{\infty}A_{n} \right)\cup \left( \bigcup _{n = k}^{\infty}B_{n}  \right)   \right)  \\ 
     &=  \bigcap_{k= 1}^{\infty}\left( C_{k}\cup D_{k} \right)  
    \end{aligned}
    \]只需要证明 \[
    \left( \bigcap_{k= 1}^{\infty}C_{k}  \right)\cup \left( \bigcap_{k= 1}^{\infty}D_{k}  \right)= \bigcap_{k= 1}^{\infty}\left( C_{k}\cup D_{k} \right)    \tag{*}
    \] 注意到 \(  C_1\supseteq C_2\supseteq \cdots   \), \(  D_1\supseteq D_2\supseteq \cdots   \), 从而 \(  C_1^{c}\subseteq C_2^{c}\subseteq \cdots   \), \(  D_1^{c}\subseteq D_2^{c}\subseteq \cdots   \),由Exercise \ref{ex:9-2-3}可得 \[
    \left( \bigcup _{n = 1}^{\infty}C_{n}^{c} \right)\cap \left( \bigcup _{n = 1}^{\infty}D_{n}^{c} \right)  = \bigcup _{n = 1}^{\infty}\left( C_{n}^{c}\cap D_{n}^{c} \right) 
    \] 其中 \[
    \left( \bigcup _{n = 1}^{\infty}C_{n}^{c} \right)=\left(  \bigcap_{n = 1}^{\infty}C_{n} \right)^{c} ,\quad  \bigcup _{n = 1}^{\infty}D_{n}^{c}=\left(  \bigcap_{n = 1}^{\infty}D_{n}   \right)^{c}  
    \] 进而 \[
    \left( \bigcup _{n = 1}^{\infty}C_{n}^{c} \right)\cap \left( \bigcup _{n = 1}^{\infty}D_{n}^{c} \right)= \left( \bigcap_{n = 1}^{\infty}C_{n}  \right)^{c}\cap \left( \bigcap_{n = 1}^{\infty}D_{n}  \right)^{c}=  \left( \left( \bigcap_{n = 1}^{\infty}C_{n}  \right) \cap  \left( \bigcap_{n = 1}^{\infty}D_{n}  \right)       \right)^{c} 
    \]此外\[
    \bigcup _{n = 1}^{\infty}\left( C_{n}^{c}\cap D_{n}^{c} \right)= \bigcup _{n = 1}^{\infty}\left( C_{n}\cup D_{n} \right)^{c}=  \left( \bigcap_{n = 1}^{\infty}\left( C_{n}\cup D_{n} \right)     \right)^{c} 
    \](*)两边的补集相等, 从而(*)成立.
    \item 令 \(  C_{k}= \bigcap_{n = k}^{\infty}A_{n}   \), \(  B_{k}= \bigcap_{ n = k}^{\infty}B_{n}   \). 则 \(  C_1\subseteq C_2\subseteq \cdots   \), \(  D_1\subseteq D_2\subseteq \cdots   \)    . \[
   \begin{aligned}
    \varliminf_{n\to\infty} (A_n \cap B_n) = \bigcup _{n = 1}^{\infty}\bigcap_{k= n}^{\infty}\left( A_{k}\cap  B_{k} \right)&=   \bigcup _{n = 1}^{\infty} \left( \left( \bigcap_{k= n}^{\infty}A_{k}  \right)\cap \left( \bigcap_{k= n}^{\infty}B_{k}  \right)  \right)\\ 
     &= \bigcup _{n = 1}^{\infty}\left( C_{n}\cap D_{n} \right) 
   \end{aligned}  
    \] \[
    \varliminf_{n\to\infty} A_n = \bigcup _{n = 1}^{\infty}C_{n},\quad  \varliminf_{n\to\infty} B_n = \bigcup _{n = 1}^{\infty}D_{n}
    \]因此只需要证明 \[
    \bigcup _{n = 1}^{\infty}\left( C_{n}\cap D_{n} \right)=  \left( \bigcup _{n = 1}^{\infty}C_{n} \right)\cap \left( \bigcup _{n = 1}^{\infty}D_{n} \right)   
    \]这由Exercise \ref{ex:9-2-3}立即得到.
    \end{enumerate}
    
\end{proof}

\begin{exercise}{}{}
   
    设$\{f_n(x)\}$以及$f(x)$都是定义在$\mathbf{R}^1$上的实值函数, 且有
    \[
        \lim_{n\to\infty}f_n(x) = f(x), \quad x \in \mathbf{R}^1,
    \]
    则对$t \in \mathbf{R}^1$, 有
    \[\begin{aligned}
        &\{x \in \mathbf{R}^1: f(x) \le t\} \\
        ={}& \bigcap_{k=1}^{\infty}\bigcup_{m=1}^{\infty}\bigcap_{n=m}^{\infty}\left\{x \in \mathbf{R}^1: f_n(x) < t + \frac{1}{k}\right\}.
    \end{aligned}\]
\end{exercise}
\begin{proof}
    记 \[
    S: = \bigcap _{k= 1}^{\infty}\bigcup _{m= 1}^{\infty}\bigcap _{n = m}^{\infty}\left\{ x\in \mathbb{R} ^{1}:  f_{n}\left( x \right)< t+ \frac{1 }{k }   \right\}
    \]
    对于固定的 \(  x  \), 若 \( \lim_{n\to \infty}f_{n}\left( x \right)=    f\left( x \right)\le t   \), 则对于任意的 \(  k\in \mathbb{Z} _{> 0}  \), 存在   \(  m \in \mathbb{Z} _{> 0}  \), 使得对于任意的 \(  n\ge m  \), 均有 \(  f_{n}\left( x \right) < t+ \frac{1 }{k }    \), 翻译成集合的语言, 即 \(  x \in S  \). 因此 \[
    \left\{ x\in \mathbb{R} ^{1}: f\left( x \right)\le t  \right\}\subseteq S
    \] 
    
    反之, 若 \(  x \in S  \),   则对于任意的 \(  k\in \mathbb{Z} _{> 0}  \), 存在   \(  m \in \mathbb{Z} _{> 0}  \), 使得对于任意的 \(  n\ge m  \), 均有 \(  f_{n}\left( x \right) < t+ \frac{1 }{k }    \). 此时, 若 \(  f\left( x \right)> t   \), 可取 \(  k_0  \)使得 \(  \frac{1 }{k_0 }< f\left( x \right)-t    \)  , 相应地 存在 \(  m_0\in \mathbb{Z} _{> 0}  \)  , 使得对于任意的 \(  n\ge m_0  \), 都有  \(  f_{n}\left( x \right)< t+ \frac{1 }{k_0 }    \). 我们得到 \[
    f\left( x \right)= \lim_{n\to \infty}f_{n}\left( x \right)\le t+ \frac{1 }{k_0 }< f\left( x \right)    
    \]   从而导出了一个矛盾. 这表明一定有 \(  f\left( x \right)\le t   \). 即 \(  x \in \left\{ x:\mathbb{R} ^{1}:f\left( x \right)\le t  \right\}  \). 因此 \[
    S\subseteq \left\{ x\in \mathbb{R} ^{1}:f\left( x \right)\le t  \right\}
    \]  
\end{proof}

\begin{exercise}{}{}
    若 $(A \setminus B) \sim (B \setminus A)$, 则 $A \sim B$, 对吗?
\end{exercise}

\begin{proof}
    正确的. 只需要证明存在 \(  A  \)到 \(  B  \)的单射, 通过调换 \(  A,B  \)的地位, 可以得到 \(  B  \)到 \(  A  \)的单射. 进而由Cantor-Bernstein-Schroeder定理, 它们是对等的. 下面来证明这件事. 由 \(  \left( A\setminus B \right)\sim \left( B\setminus A \right)    \), 可知存在单射 \(  f:A\setminus B\to B\setminus A  \). 定义 \(  g:A\to B  \) \[
    g\coloneqq\begin{cases} f\left( x \right),&x \in A\setminus B\\ 
     x,&x \in A\cap B  \end{cases} 
    \]       注意到 \(  g\left( A\setminus B \right)\cap g\left( A\cap B \right)= f\left( A\setminus B \right)\cap \left( A\cap B \right)\subseteq \left( B\setminus A \right)\cap \left( A\cap B \right)= \varnothing        \) .且 \(  g|_{A\setminus B} = f \)和 \(  g|_{A\cap B}= Id _{A\cap B}  \)均为单射. 因此 \(  g  \)本身也是一个单射.   
\end{proof}

\begin{exercise}{}{}
    若 $A \subset B$, 且 $A \sim (A \cup C)$, 试证明 $B \sim (B \cup C)$.
\end{exercise}
\begin{proof}
    首先, \[
    B \subseteq B\cup C \implies \operatorname{card}\left( B \right)\le \operatorname{card}\left( B\cup C \right)  
    \]  由于 \(  A\sim \left( A\cup C \right)   \), 存在单射 \(  f:A\cup C\to A  \). 定义 \(  g:B\cup C\to B  \)  \[
    g\left( x \right) \coloneqq \begin{cases} f\left( x \right),& x \in A\cup C\\ 
     x,& x \in B\setminus \left( A\cup C \right)      \end{cases} 
    \]  由于 \(  \operatorname{Im}g|_{A\cup C}\cap \operatorname{Im}g|_{B\setminus \left( A\cup C \right) }=   \)  \(  f\left( A\cup C \right)\cap \left( B\setminus  \left( A\cup C \right)  \right)\subseteq A\cap \left( B\setminus A \right)= \varnothing     \), 且 \(  g|_{A\cup C}  \)和 \(  g|_{B\setminus  \left( A\cup C \right) }  \)均为单射, 故 \(  g  \)是单射.  这表明 \[
    \operatorname{card}\left( B\cup C \right)\le \operatorname{card}\left( B \right)  
    \] 由Cantor-Bernstein-Schroeder定理, 我们有 \[
    \operatorname{card}\left( B\cup C \right)= \operatorname{card}\left( B \right)  
    \]即 \(  B\sim \left( B\cup C \right)   \)  
\end{proof}
\begin{exercise}{}{}
  设$E_k \subset \mathbb{R}^n (k = 1, 2, \cdots)$. 证明:
\[\begin{aligned}
\overline{\lim_{k \to \infty}} \chi_{E_k}(x) = \chi_{\overline{\lim_{k \to \infty}} E_k}(x).
\end{aligned}\]
\end{exercise}
\begin{proof}
    给定 \(  x \in \mathbb{R} ^{n}  \). 若 \(  \overline{\lim}_{k\to \infty} \chi _{E_{k}}\left( x \right)= 0   \), 则 \[
    \lim_{n\to \infty} \sup _{k\ge n}\chi _{E_{k}}\left( x \right)= 0 
    \]  这表明 存在 \(  m \in \mathbb{Z} _{> 0}  \), 使得对于任意的 \(  n\ge m  \), 都有  \[
    \sup _{k\ge n}\chi _{E_{k}}\left( x \right)< \frac{1 }{2 }  
    \]  但是 \(  \chi _{E_{k}}\left( x \right)\in \left\{ 0,1 \right\}   \), 因此只能有 \[
    \chi _{E_{k}}\left( x \right)= 0, \forall k\ge n \implies  x \in E_{k}^{c},\forall k\ge n 
    \] 以上表明 \[
    x \in  \bigcup _{m= 1}^{\infty} \bigcap _{n = m}^{\infty} E_{k}^{c}= \bigcup _{m= 1}^{\infty}\left( \bigcup _{n = m}^{\infty}E_{k} \right)^{c}= \left( \bigcap _{m= 1}^{\infty}  \bigcup _{n = m}^{\infty}E_{k} \right)^{c}= \left( \overline{\lim_{k\to \infty}}E_{k} \right)^{c}   
    \]从而 \[
    \chi _{\overline{\lim_{k\to \infty}}E_{k}}\left( x \right)= 0 
    \]

    若  \(  \overline{\lim_{k\to \infty}}\chi _{E_{k}}\left( x \right)\neq 0   \), 则 \(  1\ge   \overline{\lim_{k\to \infty}}\chi _{E_{k}}\left( x \right)> 0  \). 从而对于任意的 \(  m\in \mathbb{Z} _{> 0}   \) , 存在 \( k\ge m  \), 使得 \[
   \chi _{E_{k}}\left( x \right)> 0
    \]  同样的, 由于 \(  \chi _{E_{k}}\left( x \right)\in \left\{ 0,1 \right\}   \), 只能有 \(  \chi _{E_{k}}\left( x \right)= 1   \), 这说明 \(  \overline{\lim_{k\to \infty}}\chi _{E_{k}}\left( x \right)= 1   \) .   此外, \(  x \in E_{k}  \) . 以上表明\[
    x \in \bigcap _{m= 1}^{\infty}\bigcup _{k= m}^{\infty}E_{k}= \overline{\lim_{k\to \infty}}E_{k}
    \]即 \[
    \overline{\lim_{k\to \infty}}\chi _{E_{k}}\left( x \right)=  \chi _{\overline{\lim_{k\to \infty}}E_{k}}\left( x \right)= 1 
    \]

    这就说明了题设等式成立.
\end{proof}
\begin{exercise}{}{}
    设 $f: X \to Y, A \subset X, B \subset Y$, 试问下列等式成立吗?
\begin{enumerate}
\item[(i)] $f^{-1}(Y \setminus B) = f^{-1}(Y) \setminus f^{-1}(B)$;
\item[(ii)] $f(X \setminus A) = f(X) \setminus f(A)$.
\end{enumerate}
\end{exercise}

\begin{proof}
    \begin{enumerate}
        \item[(i)]  \[
        \begin{aligned}
        x\in f^{-1} \left( Y\setminus B \right)&\iff f\left( x \right)\in Y\setminus B\\ 
         &\iff    f\left( x \right)\in Y \text{ 且 } f\left( x \right)\not \in B\\ 
          &\iff  x \in f^{-1} \left( Y \right)\text{ 且 } x\not \in f^{-1} \left( B \right)\\ 
           &\iff x \in f^{-1} \left( Y \right)\setminus f^{-1} \left( B \right)      
        \end{aligned}
        \]故等式成立.
        \item[(ii)]等式不成立, 考虑\(  X,Y  \)非空, 且 \(  A  \)为 \(  X  \)的非空真子集.    \(  f  \)在 \(  X  \)上取常值 \(  y \in Y  \).   此时 \[
        f\left( X\setminus A \right)= \left\{ y \right\},\quad f\left( X \right)\setminus f\left( A \right)= \varnothing   
        \]
    \end{enumerate}
    
\end{proof}
\end{document}